%! AP Statistics — Unit 1, Lesson 1.5 — Warm-Up
\documentclass[10pt]{article}
\usepackage{apstats-article}
\usepackage{apstats-boxes}

%=============================================================================
\begin{document}

\pageheader{Unit 1, Lesson 1.5}{Warm-Up}
\namedateperiod

\vspace{0.16in}

\noindent The dotplot below shows the number of books read over the summer by
each of 20 students.

\vspace{0.14in}

\begin{center}
\begin{tikzpicture}[scale=1.0]
  % Axis
  \draw[thick] (0,0) -- (5.2,0);
  \foreach \x/\lbl in {0.8/1, 1.8/2, 2.8/3, 3.8/4, 4.8/5} {
    \draw (\x,-0.12) -- (\x,0.12);
    \node[below, font=\small] at (\x,-0.12) {\lbl};
  }
  \node[below, font=\small\itshape] at (2.8,-0.52) {Number of books read};
  % Dots above 1 (freq 4)
  \foreach \k in {1,...,4} \filldraw[navy] (0.8, \k*0.38) circle (0.14);
  % Dots above 2 (freq 7)
  \foreach \k in {1,...,7} \filldraw[navy] (1.8, \k*0.38) circle (0.14);
  % Dots above 3 (freq 6)
  \foreach \k in {1,...,6} \filldraw[navy] (2.8, \k*0.38) circle (0.14);
  % Dots above 4 (freq 2)
  \foreach \k in {1,...,2} \filldraw[navy] (3.8, \k*0.38) circle (0.14);
  % Dots above 5 (freq 1)
  \filldraw[navy] (4.8, 0.38) circle (0.14);
\end{tikzpicture}
\end{center}

\vspace{0.18in}

\begin{enumerate}[leftmargin=1.5em, itemsep=10pt]

  \item What is the variable being measured? Is it quantitative or categorical?
        Explain how you know.\\[6pt]
        \writeline\\[1pt]\writeline

  \item What value occurs most often in this distribution?
        How many students chose that value?\\[6pt]
        \writeline

  \item What is the \textbf{range} of values shown in the dotplot?
        (Recall: range $=$ maximum $-$ minimum.)\\[6pt]
        \writeline

  \item A classmate claims: ``A dotplot is basically the same as a bar chart for
        categorical data --- they both stack symbols above labels.''
        Name \textit{one key difference} between a dotplot for quantitative data
        and a bar chart for categorical data.\\[6pt]
        \writeline\\[1pt]\writeline

\end{enumerate}

\vspace{0.16in}

\begin{tcolorbox}[colback=greenbg, colframe=greenacc, boxrule=0.8pt, arc=2mm,
    left=4mm, right=4mm, top=3mm, bottom=3mm]
\small
\begin{itemize}[leftmargin=1.3em, itemsep=5pt]
  \item A \textbf{dotplot} places a dot above a number line for each data value ---
        it shows every individual value and reveals the overall \textit{shape} of
        the distribution.
  \item Today we add two more displays for quantitative data:
        \textbf{stemplots} and \textbf{histograms}.
  \item Each display has trade-offs: some preserve every value; others group data
        into intervals. Choosing the right display depends on the dataset size and
        the question you want to answer.
\end{itemize}
\end{tcolorbox}

\vspace{0.14in}

\noindent\textcolor{slate}{\small One question you have going into today's lesson,
or something you're not sure about yet:}\\[8pt]
\writeline

\end{document}
